\begin{table}[H]\centering\caption*{Supplementary Table S4. Robustness to attribute subsets.}
\small\begin{tabular}{lrrcc}\hline
Subset & $n$ & trend/cent. & reach [95\% CI] & rating [95\% CI]\\\hline
\quad spectacle: full (6) & 37,231 & 0.51 & 0.30 [0.29, 0.31] & 0.17 [0.15, 0.17]\\
\quad spectacle: cross-medium-validated (3) & 37,230 & 0.36 & 0.19 [0.18, 0.21] & 0.00 [-0.01, 0.01]\\
\quad spectacle: readable-structural (5) & 37,230 & 0.50 & 0.30 [0.29, 0.31] & 0.05 [0.04, 0.06]\\
\quad craft: full (5) & 37,231 & 0.71 & 0.05 [0.04, 0.06] & 0.45 [0.44, 0.45]\\
\quad craft: structural-only (3) & 37,231 & 0.92 & 0.08 [0.07, 0.09] & 0.35 [0.34, 0.36]\\
\bottomrule\end{tabular}\par\smallskip{\footnotesize\raggedright\noindent \textit{Notes.} Each row re-estimates the headline reception results on a subset of the spectacle- and craft-index attributes, on the IMDb-matched film sample ($n=37{,}231$). Reach and rating are within-decade partial correlations of the index with log vote count and average rating; the trend column is the per-century slope of standardized decade means. Bracketed intervals are 95\% confidence intervals from a nonparametric bootstrap over the matched films (600 resamples). The intervals are narrow, so the dissociation, spectacle tracking reach while craft tracks rating, is precisely estimated rather than resting on a few marginal correlations, and the attribute subsets are pre-specified definitions rather than chosen on the outcome. The full spectacle index already tilts toward reach (0.30 vs 0.17), a separation that sharpens once the index is restricted to its cross-medium-validated core (0.19 vs 0.00), and the rating effect likewise survives dropping the overtly evaluative attributes from the craft index, so the dissociation is a property of the indices rather than of every attribute definition. The per-century climb is positive across all attribute definitions.\par}\end{table}
